%! Random Sampling Methods — Unit 3, Lesson 3.2 — Lesson Plan
% GRADUAL-RELEASE plan (warm-up / guided notes and practice / debrief / close and start the
% homework) at the COMPACT budget: the whole lesson is finished in one 55-minute period.
% THERE IS NO GROUP ACTIVITY IN THIS COURSE and NO INDEPENDENT PRACTICE SET in the notes — the
% homework is the individual practice, started in class. Teacher-facing. Every box is `skillbox`
% (breakable) with a \boxguard ahead of it; no `fixedskillbox`, no tier boxes, no exit ticket.
% CED Topic 3.3 only. Systematic sampling is named in one line, with no problems on it; census
% belongs to Lesson 3.1; the kinds of bias belong to Lesson 3.3.
\documentclass[10pt]{article}
\usepackage{apstats-article}
\usepackage{apstats-boxes}

\newcommand{\UnitNumberName}{Unit 3: Collecting Data \quad}
\newcommand{\LessonNumberName}{Lesson 3.2: Random Sampling Methods}

\begin{document}
\begin{center}
    {\Large\bfseries \CourseName \\
    {\small\bfseries \UnitNumberName \LessonNumberName}}
\end{center}

\begin{tcolorbox}[colback=sky, colframe=navy, boxrule=0.9pt, arc=2mm,
    left=2mm, right=2mm, top=1.5mm, bottom=1.5mm]
    \textbf{Primary Objective:} Students will describe how to select a simple random sample
    with a random number generator or a table of random digits, identify stratified random and
    cluster samples from the description of a study, and explain in context why one sampling
    method suits a situation better than another --- in particular, that strata must be alike
    inside and clusters must each look like the population.
    \\[2pt]
    \textbf{CED alignment:} Topic 3.3 (Random Sampling and Data Collection) \quad$\cdot$\quad
    Big Idea: DAT (Data-Based Predictions, Decisions, and Conclusions) \quad$\cdot$\quad Skill
    1.C \quad$\cdot$\quad LO DAT-2.C, DAT-2.D \quad$\cdot$\quad EK DAT-2.C.1 (with and without
    replacement), DAT-2.C.2 (SRS and its mechanisms), DAT-2.C.3 (stratified), DAT-2.C.4
    (cluster), DAT-2.C.5 (systematic, named only), DAT-2.D.1 (advantages and disadvantages
    depend on the question and the population)
    \\[2pt]
    \textbf{Lesson model:} \emph{Gradual release --- I do, we do, you do --- all of it inside
    the guided notes.} There is no group activity. The notes name the vocabulary as they teach
    it and deliver the instruction row by row in one \emph{Main Ideas / Notes} table --- each
    idea a definition to complete and a grid of short problems, the first worked by you and the
    rest by the class; the last row, Guided Practice, is worked together with students holding
    the pen; \textbf{the homework is the individual practice}, and students start it in class
    while you circulate. The whole-class debrief is \emph{spoken}.
    \\[2pt]
    \textbf{Compact budget:} the notes are three pages, 14 problems. Every display is given.
    Census (Lesson 3.1) is not re-taught; the named kinds of bias (voluntary response,
    undercoverage, nonresponse, response) are held for Lesson 3.3.
\end{tcolorbox}

\renewcommand{\arraystretch}{1.4}
\boxguard
\begin{skillbox}[Learning Targets \& Key Understandings]{goldbox}
\begin{tabularx}{\linewidth}{@{}X|X@{}}
\textbf{Learning targets (student language)} & \textbf{Key understandings (the ``why'')} \\[2pt]
\hline
\vspace{-6pt}
\begin{itemize}[leftmargin=1.1em, nosep, after=\vspace{2pt}]
  \item I can describe how to choose a simple random sample with a random number generator or
        a table of random digits: label, randomize, select, skipping repeats.
  \item I can tell sampling with replacement from sampling without replacement.
  \item I can identify a stratified random sample and a cluster sample from the description of
        a study.
  \item I can explain in context why a stratified or a cluster sample is the better choice.
\end{itemize}
&
\vspace{-6pt}
\begin{itemize}[leftmargin=1.1em, nosep, after=\vspace{2pt}]
  \item In an SRS every \emph{group} of size $n$ has an equal chance to be the sample (EK
        DAT-2.C.2) --- an equal chance for each \emph{individual} is not enough: picking one
        whole grade at random gives each student a 1-in-4 chance and is not an SRS.
  \item Strata are homogeneous; take an SRS from each (DAT-2.C.3). Clusters are heterogeneous,
        each like the population; take all of a random few (DAT-2.C.4).
  \item Stratifying on a variable related to the response gives a less variable estimate;
        clustering is cheaper and more practical (DAT-2.D.1).
  \item \textbf{Target misconception:} confusing stratified with cluster --- ``the grades are
        clusters, so pick two grades and survey everyone.'' Groups alike inside on something
        related to the response are \emph{strata} (sample from every one); groups that each
        look like the whole are \emph{clusters} (all of some).
\end{itemize}
\end{tabularx}
\end{skillbox}

\boxguard
\begin{skillbox}[{Vocabulary, Concepts, \& Theorems --- taught directly in the guided notes}]{greenbox}
\begin{minipage}[t]{0.48\linewidth}
    \begin{itemize}
        \item \textbf{Simple random sample (SRS)} --- every group of $n$ individuals has an equal
              chance to be the sample; chosen by label / randomize / select
        \item \textbf{Without replacement} --- an individual can be chosen only once;
              \textbf{with replacement} --- it can be chosen again (DAT-2.C.1)
        \item \textbf{Stratified random sample} --- split into strata, groups alike inside; an
              SRS from every stratum, combined
    \end{itemize}
\end{minipage}
\hfill
\begin{minipage}[t]{0.48\linewidth}
    \begin{itemize}
        \item \textbf{Cluster sample} --- split into clusters, each like the population;
              randomly choose some and measure everyone in them
        \item \textbf{Systematic random sample} --- a random start, then every $k$th individual
              on a list (named only; no problems)
        \item \textbf{Carried forward} --- population, sample, and census (Lesson 3.1)
    \end{itemize}
\end{minipage}
\end{skillbox}

% A tabularx never splits, so guard generously ahead of it.
\boxguard[30]
\begin{skillbox}[Lesson at a Glance --- \MeetingLength]{sky}
\renewcommand{\arraystretch}{1.25}
\begin{tabularx}{\linewidth}{@{}l c X X@{}}
\textbf{Phase} & \textbf{Min} & \textbf{Students} & \textbf{Teacher} \\
\hline
Warm-Up & 5 & Lakeview's survey of 80: population and sample, names in a hat, too few seniors,
two whole advisory groups & Leave ``20 from each grade'' and ``2 whole advisory groups'' on the
board, unnamed \\
Guided Notes \& Practice & 34 & Fill the vocabulary box and the notes table: problems 1--10 row
by row, then \emph{Parking at Riverside} (11--14) together, pen in their hand & \textbf{I do}
each row's definition lines and first problem (1, 4, 7), \textbf{we do} the rest (24)
$\to$ \textbf{we do} Guided Practice (10) \\
Debrief & 8 & Pens down, papers up --- answer aloud, nothing to fill in & Two columns on the
board, \emph{strata} and \emph{clusters}; open on the almost-right problem~10 \\
Close \& Assign & 8 & \textbf{Start the homework in class}, alone and silent & Circulate on
Part~B item~3 for your formative read; preview Lesson 3.3 \\
\end{tabularx}
\end{skillbox}

\boxguard[20]
\begin{skillbox}[Warm-Up --- Activate Prior Knowledge (5 min)]{sky}
\begin{minipage}[t]{0.48\linewidth}
    \textbf{The items and what each seeds}
    \begin{itemize}
        \item \textbf{1.} Population and sample for Lakeview's survey of 80. \textbf{Spirals}
              to Lesson 3.1.
        \item \textbf{2.} Names in a hat: does everyone have the same chance? \textbf{Seeds}
              the SRS row (problems 1--3).
        \item \textbf{3.} The hat could give very few seniors, and seniors work more.
              \textbf{Seeds strata} --- the room's fix (``some from each grade'') is problem~4.
        \item \textbf{4.} Two whole advisory groups are easier to survey. \textbf{Seeds
              clusters} and the ``cheaper'' advantage (problems 5 and 8).
    \end{itemize}
\end{minipage}
\hfill
\begin{minipage}[t]{0.48\linewidth}
    \textbf{Running it}
    \begin{itemize}
        \item \textbf{Debrief items 3 and 4 aloud} and write the room's two plans on the board in
              their words --- ``some from every grade'' and ``two whole advisory groups.'' Do not
              name them. Row~2 names them ten minutes later; point back.
        \item On item~3, press for a \emph{direction}: too few seniors makes the estimate too
              low. That is the reason stratifying helps, and it returns in problem~7.
        \item Item~1 is a spiral; take 60 seconds and move on.
    \end{itemize}
\end{minipage}
\end{skillbox}

\boxguard
\begin{skillbox}[Guided Notes \& Practice --- I do, then we do (34 min)]{sky}
\textbf{This is the whole lesson.} There is no group activity, and there is no independent
practice set on this page: \textbf{the homework is the individual practice}, started in class.
The notes are one \emph{Main Ideas / Notes} table --- four rows, 14 short problems.
\textbf{In every row you fill the definition lines and work the first problem of each grid; the
class works the rest, pen in their hand.} The vocabulary box is filled \emph{as you go}.
\begin{multicols}{2}
\textbf{SRS (8 min).} Fill \emph{group}, \emph{equal}, \emph{once}, and the three steps; say
\emph{label, randomize, select} out loud. Work \textbf{problem~1} yourself --- a full
description: what is numbered, the range, skipping repeats, who is surveyed. Problem~2 is
theirs: two numbers get skipped for two different reasons (9661 is out of range, the second
0738 is a repeat). Problem~3 is the only with-replacement case today.

\textbf{Strata \& Clusters (7 min).} Fill the two definitions and read the systematic line
aloud \emph{once} --- it gets no problems. Work \textbf{problem~4} (20 from each grade:
stratified) and point at the warm-up's board. \textbf{Problem~6 is the trap}: every student
has a 1-in-4 chance, so it \emph{sounds} like an SRS. Ask: ``Could the 300 be 150 sophomores
and 150 seniors?'' No --- so not every group of 300 is possible.

\textbf{Choosing (9 min) --- the crux.} Fill \emph{differs}, \emph{varies}, \emph{expensive}.
Work \textbf{problem~7} (each grade is covered, and hours are alike within a grade).
Problems 8 and 9 are theirs; 9 is the mirror of the crux --- mixed advisory groups make
useless strata. \textbf{Problem~10 is the crux}: take a hand count on ``Jalen's plan is fine''
\emph{before} anyone explains. Then fill the callout --- \emph{every} one; \emph{all} of
\emph{some} --- and make the room say it back.

\columnbreak
\textbf{We do (about 10 min): Parking at Riverside.} A fresh context, problems 11--14, worked
entirely together. Read the table first: every building holds every apartment type --- so a
building is a small copy of the complex. Students hold the pen; you hold the questions:
\emph{``Are the households in one building alike, or a mix?''} \quad \emph{``Who owns more
cars, a one-bedroom or a three-bedroom?''} \textbf{Problem~14 is the crux transferred}: the
apartment types offered as clusters. It is the one to protect when time is short --- a student
who answers both 10 and 14 has separated \emph{alike inside} from \emph{like the whole}.

\textbf{The pen must actually change hands.} Put problem~11 on them cold (it is problem~1 again
with new numbers), take a hand count on problem~14 before anyone speaks, and make the room
defend the answer. Circulate and demand the reason:
\begin{itemize}[leftmargin=1.1em, nosep]
  \item \textbf{Problem~12:} ``Why is one trip to 5 buildings cheaper than 60 doors?''
  \item \textbf{Problem~13:} ``What do apartments of one type have in common?''
  \item \textbf{Problem~14:} ``If you got only the one-bedrooms, would the mean be too high or
        too low?''
\end{itemize}
While circulating, pick the papers to put up in the debrief --- one \emph{almost}-right
problem~10 (``wrong, those are strata,'' with no reason) and one that names what a cluster
must look like. \textbf{Your formative read comes twice}: here, and again in the last eight
minutes as students start the homework.
\end{multicols}
\end{skillbox}

\boxguard
\begin{skillbox}[Debrief --- whole class (8 min)]{sky}
Pens down, papers up. \textbf{This debrief is spoken} --- nothing on the page to fill in.
\begin{multicols}{2}
\begin{enumerate}[leftmargin=1.3em, nosep, itemsep=3pt]
  \item \textbf{Open on the almost-right problem~10}: ``Jalen is wrong --- grades are strata.''
        Right label, no reason. Ask the room for the reason: grades differ on hours worked, so
        two whole grades could leave out the seniors and the estimate is too low. Correct it
        gently and in public: the label follows from \emph{alike inside} versus \emph{like the
        whole}.
  \item Two columns on the board: \emph{Strata --- sample from every one}: Lakeview's grades,
        Riverside's apartment types. \emph{Clusters --- all of some}: advisory groups,
        buildings. That board is the lesson.
  \item Take problem~6 back: why is ``everyone had a 1-in-4 chance'' not enough for an SRS?
\end{enumerate}
\columnbreak
\textbf{The sentence to leave on every desk:} \emph{Strata are alike inside --- sample from
every one. Clusters each look like the whole --- take all of some.} Make a student say it, not
you.

\textbf{If the debrief runs short}, cut item~3 --- item~1 is the only one that cannot be cut.
\textbf{Do not let the debrief eat the last eight minutes} --- the homework start is your
second formative read.
\end{multicols}
\end{skillbox}

\boxguard
\begin{skillbox}[AP Practice --- extra credit, scored 1 to 5]{goldbox}
Two multiple-choice items --- \textbf{MC~1} names the method (8 of 90 counties at random, every
farm in them: a cluster sample, \textbf{B}; the distractor C calls the counties strata), and
\textbf{MC~2 is the crux in AP clothing}: why stratify textbook spending by college
(\textbf{C}, less variability), with (D) ``cheaper, only a few colleges visited'' and (E)
``survey everyone in each chosen college'' both cluster ideas wearing a stratified label ---
then one free-response set on a full flight: \textbf{(a)} describe an SRS of 24 with a random
number generator, \textbf{(b)} a stratified sample by window / middle / aisle and why it beats
an SRS, and \textbf{(c)}, \textbf{the spiral}, the population and the sample (Lesson 3.1).

\textbf{Extra credit, scored 1 to 5, turned in with the homework.} Go over it at the start of
Lesson 3.3. Hold the AP standard out loud: \emph{in context or it does not count} --- a part (b)
that says ``stratified is more accurate'' without saying comfort is alike within a seat type
has not earned the point, and a part (a) without ``skip repeats'' or without saying who is
surveyed is partial.
\end{skillbox}

\boxguard
\begin{skillbox}[Homework --- scored 1 to 5, due the first class after two study halls]{goldbox}
Six items in a third context: Mercy Valley Hospital's 720 nurses, on three shifts and twelve
floors, each floor staffed from every shift; the question is hours of sleep, and night-shift
nurses sleep less. \textbf{Part~A}: \textbf{A1} name the floor plan (cluster); \textbf{A2} name
the shift plan (stratified); \textbf{A3} an AP-style multiple choice --- which plan is most
precise (C, stratify by shift) --- with a required justification. \textbf{Part~B}: \textbf{B1}
describe an SRS of 36 with a random number generator; \textbf{B2} a practical advantage of the
floor (cluster) plan; \textbf{B3} the crux restated --- a supervisor wants the shifts as
clusters.

\textbf{Start it in the last eight minutes of the period, in class and alone.}
\textbf{Scoring:} 10 points --- A1, A2, A3's letter, and A3's justification 1 each; B1, B2, B3
2 each. \textbf{B3 is the one to read closely}: a student who accepts the supervisor's plan has
kept the target misconception, and Lessons 3.4--3.5 (blocking is stratifying for experiments)
build on its absence. The spiralbox previews Lesson 3.3.

\textbf{Paper homework is the default for this course} --- assign this page unless you have
decided otherwise. \textbf{DeltaMath override (occasional):} on the days you assign DeltaMath
instead, it replaces this page, and this page stays in the packet as extra practice. Never
assign both.
\end{skillbox}

\boxguard
\begin{skillbox}[Watch For (while circulating)]{redbox}
\begin{itemize}[leftmargin=1.2em, nosep]
  \item \textbf{Strata and clusters swapped} (notes~9, 10, 14; HW~B3; AP~MC1, MC2 --- the target
        misconception). Probe: \emph{``Are the people inside one group alike, or a mix? If
        alike, sample from every group. If each group looks like the whole, take all of a
        few.''}
  \item \textbf{``Everyone had the same chance, so it is an SRS''} (notes~6). Probe:
        \emph{``Could the sample be half sophomores and half seniors? Then not every group was
        possible.''}
  \item \textbf{An SRS description with no mechanics} (notes~1 and 11, HW~B1, AP~(a)): ``pick 80
        at random.'' Probe: \emph{``What gets numbered? What range? What about a repeat? Who is
        surveyed?''}
  \item \textbf{``Stratified is better'' with no reason} (notes~7 and 13, HW~A3, AP~(b)).
        Probe: \emph{``Better how --- and what is alike inside each stratum?''}
\end{itemize}
Cold-call prompts: \emph{``Name a group at our school that would make a good stratum for
`hours of sleep' --- and one that would make a good cluster.''} \quad \emph{``Why does a cluster
sample save money?''}
\end{skillbox}

\boxguard
\begin{skillbox}[Close \& Assign (8 min)]{goldbox}
\textbf{Students start the homework here, in class, alone and silent --- this is the individual
practice, and the first minutes of it are your best formative read.} Hand the launch first ---
six items on paper, scored, due the first class after two study halls --- and point out that
the AP Practice page is extra credit, turned in with the homework. Circulate on \textbf{Part~B
item~3}. Sort what you see into three piles as you walk --- rejected the supervisor's plan and
said shifts should be strata, with the reason (ready); rejected it with no reason, or only
``that's stratified'' (ready, needs the language); accepted it because the shift is picked at
random (the misconception survived) --- and open Lesson 3.3 by returning to the
strata / clusters board if that third pile ran past a third of the room. Then close on what
changed: \emph{``You walked in with one kind of random sample. You walk out asking whether the
groups are alike inside or each like the whole.''} \textbf{Preview:} Lesson 3.3, \emph{Problems
with Sampling} --- even a random method can go wrong when part of the population is left out
or people will not answer.
\end{skillbox}

% ── Teacher notes — one per component, in packet order (four: there is no activity) ──
\begin{teachernote}[Warm-Up]
Item~3 is the seed for strata and item~4 the seed for clusters. Give five minutes total and
spend at most a minute on item~1. Write the room's two plans on the board in their words and
\emph{do not name them}: the notes name them, and the point-back is the payoff. On item~2,
accept ``yes, if the names are mixed well''; the deeper point (equal chance for every
\emph{group}) comes in problem~6.
\end{teachernote}

\begin{teachernote}[Guided Notes \& Practice]
\textbf{Pace it 24 / 10, then 8 for the debrief, then 8 for the homework start.} Four rows, 14
problems. \textbf{You work 1, 4, and 7}; the class works the rest. The SRS row must move ---
eight minutes --- because the time belongs to the Choosing row, which carries the crux.
\textbf{Problem~6 is the trap} (a 1-in-4 chance for every student is not an SRS);
\textbf{problem~10 is the crux} (the grades offered as clusters). Take a hand count on 10
before any explanation; the callout sentence is the resolution. The systematic line is read
once and left --- do not build problems on it. \textbf{Parking at Riverside is where the pen
changes hands}; problem~14 is the crux transferred and the one to protect. If you only get
through 11, 13, and 14, that is the right trade. \textbf{Your second formative read is the
homework start}: Part~B item~3, three piles. \textbf{Debrief (8 min), spoken}, opening on the
almost-right problem~10. \textbf{Do not let the debrief eat the last eight minutes} --- the
homework start is not optional padding.
\end{teachernote}

\begin{teachernote}[AP Practice]
\textbf{This page is extra credit, scored 1 to 5} --- the graded practice is the homework
behind it. \textbf{MC~2 is the one worth reading closely}: (D) and (E) describe a cluster
sample's features; if more than a handful choose either, return to the strata / clusters board
before Lesson 3.3. On the free response, part (a) needs the labels, the range, skipping
repeats, and who is surveyed; part (b) needs the three strata, a sample from \emph{each}, and
the reason in context (comfort is alike within a seat type); part (c) is a Lesson 3.1 spiral ---
the population is the 240 passengers on this flight, not ``airline passengers.''
\end{teachernote}

\begin{teachernote}[Homework]
\textbf{Scored 1 to 5 (10 points), due the first class after two study halls.} The eight
minutes at the end of the period are a \emph{start}: put students on A1, A2, and B3 first while
you can still catch a wrong start. \textbf{B3 is the item to read first when scoring} --- a
student who accepts the supervisor's plan has kept the target misconception. On A1 and A2,
accept ``cluster'' and ``stratified'' with or without ``sample.'' On A3, (C) with no
justification earns 1 of 2. On B1, both the labeling and the skip-repeats / who-is-surveyed
step are needed for 2 points. On B2, the advantage must be practical (cost, time, travel), not
``more accurate.'' \textbf{Paper is the default.} On the occasions you assign DeltaMath in its
place, tell students this page is now optional practice. Do not assign both.
\end{teachernote}
\end{document}
